% !TEX program = xelatex
\documentclass[12pt,a4paper]{ctexart}

\usepackage{fontspec}
\setCJKmainfont{SimSun}[BoldFont=SimHei,ItalicFont=KaiTi]
\setCJKsansfont{SimHei}
\setCJKmonofont{FangSong}

\usepackage[top=28mm,bottom=28mm,left=28mm,right=28mm]{geometry}
\usepackage{setspace}
\usepackage{fancyhdr}
\usepackage{hyperref}
\usepackage{perpage}
\usepackage{xcolor}
\MakePerPage{footnote}

\hypersetup{colorlinks=true, linkcolor=black, urlcolor=blue!50!black}

\setlength{\parindent}{2em}
\setlength{\parskip}{0pt}
\linespread{1.35}

\pagestyle{fancy}
\fancyhf{}
\fancyfoot[C]{\thepage}
\renewcommand{\headrulewidth}{0pt}

\begin{document}

\begin{center}
    {\zihao{2}\bfseries 《新民主主义论》读书报告\par}
    \vspace{1.2em}
    {\zihao{-4}
    课程名称：毛泽东思想和中国特色社会主义理论体系概论\par}
    \vspace{0.4em}
    {\zihao{-4}
    阅读篇目：毛泽东《新民主主义论》\par}
    \vspace{0.4em}
    {\zihao{-4}
    姓名：\underline{\hspace{4.5cm}} \qquad
    学号：\underline{\hspace{4.5cm}}\par}
    \vspace{0.4em}
    {\zihao{-4}
    提交日期：2026年4月18日\par}
\end{center}

\vspace{1.5em}

在本学期毛概课程的推荐阅读中，我选择《新民主主义论》作为读书报告对象。吸引我的，是它开篇提出的“中国向何处去”。在1940年的历史条件下，这并非抽象追问，而是关系民族前途与革命道路的现实问题。重读原文后我感到，它最重要的价值不只是提出“新民主主义”概念，而是通过分析中国社会性质和革命阶段，把时代问题转化为一条较为清晰的实践道路。\footnote{毛泽东：《新民主主义论》，《毛泽东选集》第2卷，北京：人民出版社，1991年，第662—706页。}

\section*{一、从时代追问进入理论论证}

这篇文章的问题意识非常鲜明。毛泽东并不是先设定一套抽象理论，再让中国现实去适应它，而是从民族危机、社会矛盾和革命处境出发，逐步推出中国需要的新政治、新经济和新文化。文章从“中国向何处去”切入，随后转入对中国社会性质的分析，这种写法说明：理论不是悬浮的，而是直接服务于现实判断和实践选择的。

过去我常把“半殖民地半封建社会”“两步走”“新民主主义的政治经济文化”当成分散知识点记忆，读原文后才意识到它们由同一逻辑贯穿：先认识国情，再决定道路。正因为中国既不同于西方资本主义国家，也不能脱离历史条件直接进入社会主义，所以中国革命既不能照搬旧式民主革命模式，也不能空谈“直接社会主义化”。\footnote{毛泽东：《新民主主义论》，《毛泽东选集》第2卷，北京：人民出版社，1991年，第664—667页。}

\section*{二、“新”字背后的历史含义}

我认为，理解《新民主主义论》的关键之一，在于理解“新”字。它之所以“新”，首先因为中国革命处在帝国主义时代，已不能纳入旧世界资产阶级民主革命的轨道；其次因为领导力量变了，真正能够把民族独立和社会变革统一起来的是无产阶级及其政党；再次因为它的前途不是停在资本主义，而是为社会主义创造条件。

由此可见，新民主主义不是旧民主主义的重复，而是在新的世界条件下对中国革命历史位置的重新界定。它既要完成民族独立和反封建任务，又把社会主义作为未来方向，因此具有鲜明的过渡性和创造性。\footnote{毛泽东：《新民主主义论》，《毛泽东选集》第2卷，北京：人民出版社，1991年，第666—667页。}

这一点也解释了为什么文章既强调民族独立，又坚持社会主义方向。若只看到前者，容易把它理解成一般的民族革命；若只看到后者，又容易忽略当时中国首先必须解决的现实任务。新民主主义的重要意义，恰恰在于把阶段性任务和长远目标统一起来。

\section*{三、新中国的构想是一整套重建方案}

《新民主主义论》给我的另一个重要印象，是它讨论的并非单纯的政权更替，而是一整套新中国方案。文章把新政治、新经济、新文化连在一起，说明革命目标不仅是“夺取政权”，更是重建社会秩序。

其中我尤其关注“民族的、科学的、大众的文化”。这句话在复习时常被当作固定表述，但放回原文，它实际上与反帝反封建、启发民众和改造旧文化紧密相连。也就是说，新民主主义既回答“谁掌权”，也回答“建设什么样的社会”和“形成什么样的精神面貌”。这正体现了《新民主主义论》的总体性思维。

更重要的是，文章把文化问题放在和政治、经济同等重要的位置，本身就说明革命不是单靠政权更替就能完成的。旧制度要真正被克服，还需要旧观念、旧习惯和旧文化权威的松动。没有这一层面的重建，新国家也难以形成稳定的社会基础。

\section*{四、和《中国革命和中国共产党》对读后的理解}

把《新民主主义论》和《中国革命和中国共产党》对读，会更清楚地看到毛泽东的思考方式。前者着重回答“应建立什么样的新中国”，后者系统分析革命对象、任务、动力和性质。两篇文章放在一起看，可以发现毛泽东始终围绕中国社会的特殊性质、阶级结构和历史任务来思考道路，而不是从抽象教条出发。这对我最大的启发是：真正有效的理论不是机械搬用现成结论，而是在把握现实条件基础上的创造性概括。\footnote{毛泽东：《中国革命和中国共产党》，《毛泽东选集》第2卷，北京：人民出版社，1991年，第636—650页。}

对今天的学习者而言，这种方法论意义并不亚于文章中的具体判断。阅读《新民主主义论》后，我更清楚地认识到：面对复杂现实，先认识对象、再判断阶段、最后选择道路，本身就是一种重要的理论能力。

例如，毛泽东反复强调先分析中国社会性质，再谈革命对象、任务和前途，这种层层推进的论证方式，使我看到理论文章也可以具有很强的现实针对性。它不是只给出结论，而是展示结论怎样从历史条件中一步步推导出来。

\section*{结语}

总体来看，《新民主主义论》的重要性，不仅在于它系统阐述了新民主主义革命的基本思想，更在于它把“中国向何处去”的时代追问转化为逻辑完整、现实明确的革命方案。通过这次阅读，我更能理解“新”不是口号式的新，而是在历史条件、领导力量和发展方向变化之后，对中国革命道路作出的创造性回答。对我来说，这种把阶段任务与长远方向统一起来的能力，正是这篇文章至今仍值得重读的原因。

\end{document}
